% !TEX program = xelatex
\documentclass[12pt,a4paper]{ctexart}

% ============ 页面与字体 ============
\usepackage{geometry}
\geometry{left=2.5cm,right=2.5cm,top=2.5cm,bottom=2.5cm}

% ============ 数学与符号 ============
\usepackage{amsmath,amssymb,amsthm}
\usepackage{mathrsfs}

% ============ 超链接与引用 ============
\usepackage{hyperref}
\hypersetup{
    colorlinks=true,
    linkcolor=blue,
    citecolor=blue,
    urlcolor=blue,
}

% ============ 定理类环境 ============
\theoremstyle{definition}
\newtheorem{definition}{定义}[section]
\newtheorem{example}{例}[section]

\theoremstyle{plain}
\newtheorem{theorem}{定理}[section]
\newtheorem{lemma}[theorem]{引理}
\newtheorem{proposition}[theorem]{命题}
\newtheorem{corollary}[theorem]{推论}
\newtheorem{conjecture}{猜想}[section]
\newtheorem{remark}{注记}[section]

\usepackage{cleveref}
\crefname{theorem}{定理}{定理}
\crefname{lemma}{引理}{引理}
\crefname{proposition}{命题}{命题}
\crefname{corollary}{推论}{推论}
\crefname{definition}{定义}{定义}
\crefname{example}{例}{例}
\crefname{remark}{注记}{注记}
\crefname{equation}{式}{式}
\crefdefaultlabelformat{#2#1#3}

% ============ 其他宏包 ============
\usepackage{graphicx}
\usepackage{enumitem}
\usepackage{booktabs}

% ============ 自定义命令 ============
\newcommand{\R}{\mathbb{R}}
\newcommand{\N}{\mathbb{N}}
\newcommand{\Q}{\mathbb{Q}}
\newcommand{\Z}{\mathbb{Z}}
\newcommand{\abs}[1]{\left|#1\right|}
\newcommand{\norm}[1]{\left\|#1\right\|}

\title{\textbf{线性逼近算子的收敛性与逼近阶——从Korovkin定理到Jackson型估计}}
\author{矜持}
\date{\today}

\begin{document}
\maketitle
\tableofcontents

% 包含各章节
\section{引言}

函数逼近论研究一个基本问题：如何用简单函数（多项式、三角多项式等）去近似复杂的连续函数？这一问题贯穿数学分析的诸多分支，Weierstrass 的两个逼近定理从原则上肯定了连续函数可用多项式或三角多项式一致逼近，但并未给出具体的构造方法。

逼近论的发展由此分化为两个相互关联的研究方向。

其一是\textbf{最佳逼近理论}。给定逼近工具的“次数”约束，是否存在一个最优的逼近元，使得逼近误差达到理论上的下确界？这个最优逼近元具有怎样的特征？切比雪夫交错定理对这一问题给出了精确的刻画——误差函数在足够多的点上达到最大值且符号交替。最佳逼近理论确立了一个不可逾越的“最优基准”，但除了少数特例，直接求出最佳逼近多项式往往极其困难。

其二是\textbf{构造性逼近理论}。能否设计结构简单、便于计算的线性算子序列，使之能一致逼近任意连续函数？Bernstein 多项式、Fejér 算子等都是这一思路的经典产物。这些算子虽然未必是最优的，但胜在形式明确、可直接计算。

由此自然产生两个核心问题：
\begin{enumerate}
    \item 对于给定的线性算子序列，如何简单有效地判定它是否对一切连续函数一致收敛？
    \item 构造出来的线性算子，其逼近效果能有多好——或者说，它们的逼近误差能否与最佳逼近误差相比拟？
\end{enumerate}

本文的主线即是围绕这两个问题展开。Korovkin 定理给出了第一个问题的答案：对于线性正算子，只需验证三个简单的“测试函数”（$1, x, x^2$ 或 $1, \cos t, \sin t$）的逼近行为，就能保证对一切连续函数的一致收敛性。这一判别法极为简洁，它将众多经典算子的收敛性证明统一到了同一个框架之下。

第二个问题则需要将构造性逼近拉回到最佳逼近的参照系中来回答。Lebesgue 不等式提供了连接二者的桥梁：任何线性算子的逼近误差，可以被最佳逼近误差乘以该算子的 Lebesgue 常数所控制。而 Jackson 定理进一步给出了最佳逼近误差的定量估计——它由函数的光滑性（连续模）所决定。综合二者，即可对具体算子的逼近阶给出精确刻画。

本文的组织结构如下：第二章介绍最佳逼近的存在性、唯一性与切比雪夫交错定理，确立逼近的“最优基准”；第三章阐述 Korovkin 定理及其在 Bernstein、Fejér 等经典算子上的应用，展示构造性逼近的统一收敛性判据；第四章引入 Lebesgue 不等式，将线性算子的误差纳入最佳逼近的框架中进行分析；第五章利用 Jackson 定理与 Zygmund 类，对逼近阶进行定量刻画。最后给出总结与展望。

\section{最佳逼近}

逼近论研究两个基本问题。一是“存在性问题”：给定函数类，能否用某类简单函数一致逼近？Weierstrass 定理对此给出了肯定的回答。二是“最优性问题”：在次数受限的逼近工具中，理论上的最小误差是多少？达到该最小误差的逼近元具有什么特征？本章讨论第二个问题，建立最佳逼近的存在性、唯一性与切比雪夫特征刻画。这些结果将为后续 Jackson 定理提供基础，也为衡量构造性算子的逼近品质提供“最优”的参照基准。

\subsection{准备工作与基础}
先约定一些记号。
\begin{definition}
\begin{enumerate}
    \item $C[a,b]$ 表示闭区间 $[a,b]$ 上的连续函数全体；
    \item $C_{2\pi}$ 表示 $\R$ 上周期为 $2\pi$ 的连续函数全体；
    \item $P_n$ 为次数不超过 $n$ 的实多项式全体
    $\{a_nx^n+\cdots+a_1x+a_0: a_i\in\R\}$；
    \item $T_n$ 为次数不超过 $n$ 的三角多项式全体
    $\{a_n\cos nx+b_n\sin nx+\cdots +a_1\cos x+b_1\sin x +a_0: a_i,b_i\in\R\}$；
    \item 约定用 $p$ 表示多项式，$t$ 表示三角多项式；
    \item $\|f\| = \max_{x\in D} |f(x)|$，其中 $D$ 为函数的定义域。
\end{enumerate}
\end{definition}

Weierstrass 第一逼近定理断言，$[a,b]$ 上的连续函数可用多项式一致逼近；第二逼近定理断言，$2\pi$ 周期连续函数可用三角多项式一致逼近。

\subsection{误差尺度的衡量}
用次数不超过 $n$ 的多项式（或三角多项式）逼近 $f$，自然关心能达到的最小误差。定义
\[
e(f,P_n) = \inf\{\|f-p\| : p\in P_n\},
\]
以及 $e(f,T_n)$ 类似。易见对任意 $n$，存在子列使得该误差趋于零，但收敛速度可能很慢（例如 $f(x)=|x|$ 在 $0$ 附近的逼近）。因此需要在限定次数的条件下寻找“最佳逼近”。

\subsection{最佳逼近的存在性}
\begin{definition}
称 $p\in P_n$ 为 $f$ 在 $P_n$ 中的\textbf{最佳逼近}，若 $\|f-p\| = e(f,P_n)$。类似定义三角多项式的最佳逼近。
\end{definition}

\begin{theorem}
设 $f\in C[a,b]$，则存在 $p\in P_n$ 使 $\|f-p\| = e(f,P_n)$。
\end{theorem}
\begin{proof}
先给出一个引理。
\begin{lemma}
设 $\{p_m\}$ 是次数不超过 $n$ 的有界多项式列，则存在收敛子列 $p_{m_k}$ 满足 $\lim_{k\to\infty} p_{m_k} = p^*$。
\end{lemma}
引理可由反复使用 Bolzano–Weierstrass 定理证明。

下记 $e_n(f)=e(f,P_n)$。对每个 $m\ge1$，存在 $p_m\in P_n$ 使得
\[
\|f-p_m\|  \leqslant  e_n(f) + \frac1m  \leqslant  e_n(f)+1.
\]
因为 $e_n(f) \leqslant  \|f-0\| = \|f\|$，故
\[
\|p_m\|  \leqslant  \|f\| + \|p_m-f\|  \leqslant \|f\| + e_n(f)+1  \leqslant 2\|f\|+1.
\]
即 $\{p_m\}$ 有界，从而存在子列 $\{p_{m_k}\}$ 收敛到某个 $p^*\in P_n$。
下证 $p^*$ 即为最佳逼近。因为
\[
\|f-p_{m_k}\|  \leqslant e_n(f) + \frac1{m_k}  \leqslant e_n(f) + \frac1k,
\]
令 $k\to\infty$ 得 $\|f-p^*\| \leqslant e_n(f)$。又因 $e_n(f)$ 是下确界，必有 $\|f-p^*\| = e_n(f)$。
\end{proof}

类似可证
\begin{theorem}
设 $f\in C_{2\pi}$，则存在 $t\in T_n$ 使 $\|f-t\| = e(f,T_n)$。
\end{theorem}

\subsection{最佳逼近的特征刻画——切比雪夫交错定理}
\begin{definition}[交错偏离点]
设 $f\in C[a,b]$，$p\in P_n$，$r(x)=f(x)-p(x)$，
若存在点列
\[
a \leqslant  x_1 < x_2 < \dots < x_N  \leqslant  b,
\]
满足：
\begin{enumerate}
    \item 对每个 $j$，$|r(x_j)| = \|r\| $；
    \item 对每个 $j=1,\dots,N-1$，$r(x_j) = -r(x_{j+1})$，
\end{enumerate}
则称 $\{x_1,\dots,x_N\}$ 为 $r$ 的一组\textbf{交错点}，每个 $x_j$ 称为\textbf{交错偏离点}。
\end{definition}

\begin{theorem}[切比雪夫最佳逼近定理]
设 $f\in C[a,b]$，则 $p\in P_n$ 是 $f$ 的最佳逼近当且仅当误差函数 $r(x)=f(x)-p(x)$ 存在 $n+2$ 个交错偏离点。\\
设 $f\in C_{2\pi}$，则 $t\in T_n$ 是 $f$ 的最佳逼近当且仅当误差函数 $r(x)=f(x)-t(x)$ 存在 $2n+2$ 个交错偏离点。
\end{theorem}
% \begin{proof}
仅对多项式情形给出证明，三角情形类似。

\textbf{充分性}：设 $x_1<\cdots<x_{n+2}$ 是 $p$ 的一组交错偏离点，即
\[
|f(x_j)-p(x_j)| = e_n(f),\quad j=1,\dots,n+2,
\]
且相邻点误差符号交替。不妨设 $f(x_j)-p(x_j) = (-1)^{j-1}e_n(f)$。
若 $p$ 不是最佳逼近，则存在 $q\in P_n$ 使 $\|f-q\| < \|f-p\|$。
考虑 $d(x)=p(x)-q(x)\in P_n$。在 $x_j$ 处，
\[
d(x_j) = (p(x_j)-f(x_j)) - (q(x_j)-f(x_j)).
\]
因 $\|f-q\|<\|f-p\|$，故 $d(x_j)$ 的符号与 $p(x_j)-f(x_j)$ 相同。于是 $d$ 在 $x_j$ 处依次变号，从而 $d$ 至少有 $n+1$ 个零点，但 $d$ 是 $n$ 次多项式，于是 $d$ 至少有$ n+1$ 个零点。
但 $d$ 为次数不超过 $n$ 的多项式，
除非 $d\equiv0$，
否则不可能有 $n+1$ 个不同零点。
故 $d\equiv0$，即 $p\equiv q$，
与假设矛盾。

\textbf{必要性}：
\begin{enumerate}
    \item \textbf{交错偏离点有限且小于 $n+2$} 

若交错偏离点个数 $N < n+2$，但 $p$ 仍是最佳逼近。记 $r(x)=f(x)-p(x)$，$M=\|r\|$。设
\[
a<x_1<\cdots<x_N<b
\]
为全体交错偏离点。

对每个 $i=1,\dots,N-1$，
由连续性可知存在
\[
\xi_i\in(x_i,x_{i+1})
\]
使得
\[
r(\xi_i)=0.
\]

再记
\[
\xi_0=a,\qquad \xi_N=b.
\]将 $[a,b]$ 用 $$\xi_1,\dots,\xi_{N-1}$$ 分割成若干子区间，则 $r$ 在每个子区间上满足
\[
-M+\alpha < r(x)  \leqslant M \quad\text{或}\quad -M  \leqslant r(x) < M-\alpha.
\]
因为我们可以假设在区间 $[a, \xi_1]$ 上有 $r(x)  \leqslant 0$。因为 $\|r\| = M$，从而 $|r(x)|  \leqslant  M$，从而此时 $r(x) \geqslant -M$ 是成立的。

由闭区间上连续函数的性质，存在 $0 < \alpha_1 < M $ 使得 $-M < r(x) < M -\alpha_1$ 在 $[a, \xi_1]$ 上成立。同理在每个子区间上都存在 $ 0 < \alpha_i < M $ 使得上述不等式成立。取 $\alpha = \min\{\alpha_i\}$，则在每个子区间上都有
\[-M   \leqslant  r(x) < M - \alpha.
\]或者
\[-M + \alpha < r(x)  \leqslant  M \]成立。

由于 $r$ 在每个子区间上不变号，
且 $\phi$ 在穿过每个 $\xi_i$ 后也恰变号一次，
故 $r$ 与 $\phi$ 在每个子区间上同号。

令 $\phi(x) = \prod_{i=1}^{N-1}(x-\xi_i)\in P_{N-1}\subseteq P_n$。构造 $q(x)=p(x)+c\phi(x)$。我们通过选取合适的 $c$ 去导出 $q(x)$ 是比 $p(x)$ 更好的逼近，从而导致矛盾。

不妨设在 $[a, \xi_1]$ 上 $r(x) \leqslant  0$ 并且 $\phi(x)$ 在 $[a, \xi_1]$ 上满足 $\phi(x)  \leqslant  0$（这只是符号上的问题）。
记 $K = \|\phi\|$.

此时 $-M  \leqslant  r(x) < M-\alpha $  注意到 $\phi(x)$ 与 $r(x)$ 在 $[a, \xi_1]$ 上符号相同。因为
$$
\begin{aligned}
f(x) - q(x) &= f(x) - p(x) - c\phi(x)  \\
              &= r(x) - c\phi(x) \\
              &= r(x) + c|\phi(x)| \\
\end{aligned}$$
这里由于 $r(x)  \leqslant  0$ 并且 $\phi(x)  \leqslant  0$，所以 
当 $c$ 满足 $ 0 < c| \phi(x) |  \leqslant  |c\cdot K| < \frac{\alpha}{2}$ 并且 $c>0$ 时
有
$$
\begin{aligned}
-M  < r(x) - c\phi(x) < M -\alpha + \frac{\alpha}{2} \\ 
\end{aligned}
$$
同理，在区间 $[\xi_1, \xi_2]$ 上我们能够利用
$$
\begin{aligned}
 -M +\alpha < r(x)  \leqslant  M
\end{aligned}
$$
和
$$
\begin{aligned}
f(x) - q(x) &= f(x) - p(x) - c\phi(x)  \\
              &= r(x) - c\phi(x) \\
              &= r(x) - c|\phi(x)| \\
\end{aligned}
$$
得到
$$
\begin{aligned}
-M + \alpha - \frac{\alpha}{2} < r(x) - c\phi(x) < M 
\end{aligned}$$
注意到上述过程与子区间的选择无关，因此在每个子区间上都满足上述不等式。
从而有 $$\|f-q\| < M = \|f-p\|$$，并且根据 $q$ 的构造可以得知它是不超过$n$ 次的多项式，
于是与 $\|f-q\| < \|f-p\|$矛盾。

\item \textbf{交错偏离点有限且大于等于 $n+2$} 

由于交错偏离点按符号交替出现，
故可从中依次选取
$n+2$
个点构成一组交错点。
由充分性部分，
$p$
即为最佳逼近。

\item \textbf{交错偏离点有无穷个}记若交错偏离点有无穷多个。记
$$
E_+=\{x\in[a,b]:r(x)=M\},\qquad
E_-=\{x\in[a,b]:r(x)=-M\}.
$$

由于 $r$ 连续，而 $\{M\},\{-M\}$ 为闭集，因此
$$
E_+=r^{-1}(\{M\}),\qquad
E_-=r^{-1}(\{-M\})
$$
均为闭集，从而为紧集。

由于交错偏离点无穷多个，因此在区间 $[a,b]$ 中，$E_+$ 与 $E_-$ 必无限次交替出现。

先取
$$
x_1=\min(E_+\cup E_-).
$$
不妨设
$$
r(x_1)=M.
$$

由于交错偏离点无穷多个，因此集合
$$
\{x>x_1:x\in E_-\}
$$
非空。又因为 $E_-$ 为闭集，所以该集合在 $[a,b]$ 中取到最小元。定义
$$
x_2=\min\{x>x_1:x\in E_-\}.
$$
于是
$$
x_1<x_2,\qquad r(x_2)=-M.
$$

同理，由于交错偏离点无穷多个，集合
$$
\{x>x_2:x\in E_+\}
$$
非空，并且取到最小元。定义
$$
x_3=\min\{x>x_2:x\in E_+\}.
$$
于是
$$
x_2<x_3,\qquad r(x_3)=M.
$$

继续递归构造：
若已定义
$$
x_k,
$$
则定义
$$
x_{k+1}
=
\min\{x>x_k:x\in E_{\sigma_{k+1}}\},
$$
其中
$$
\sigma_{k+1}=
\begin{cases}
+,& r(x_k)=-M,\\
-,& r(x_k)=M.
\end{cases}
$$

由于交错偏离点无穷多个，上述过程不会在有限步停止。

因此可得到
$$
x_1<x_2<\cdots<x_{n+2}
$$
满足
$$
r(x_i)=(-1)^iM.
$$

故已经存在不少于 $n+2$ 个交错偏离点。
\end{enumerate}
\end{proof}

证明比较繁琐，详细见附录内容。
\subsection{最佳逼近的唯一性}
\begin{theorem}
最佳逼近是唯一的：对 $f\in C[a,b]$，存在唯一的 $p\in P_n$ 使得 $\|f-p\| = e_n(f)$；对 $f\in C_{2\pi}$，存在唯一的 $t\in T_n$ 使得 $\|f-t\| = e(f,T_n)$。
\end{theorem}

\begin{proof}
若不然，设 $p,q\in P_n$ 均为最佳逼近，则 $\|f-p\| = \|f-q\| = e_n(f)$。由三角不等式，
\[
\left\|f - \frac{p+q}{2}\right\|  \leqslant \frac12\|f-p\| + \frac12\|f-q\| = e_n(f),
\]
故 $\frac{p+q}{2}$ 亦为最佳逼近。由切比雪夫交错定理，存在 $n+2$ 个交错偏离点 $x_1<\cdots<x_{n+2}$。记 $d_p = f-p$, $d_q = f-q$，则
\[
\left|f(x_1) - \frac{p(x_1)+q(x_1)}{2}\right| = \frac{|d_p(x_1)+d_q(x_1)|}{2} = e_n(f).
\]
因 $|d_p|,|d_q| \leqslant e_n(f)$，上式仅当 $d_p(x_1)=d_q(x_1)$ 时成立，从而 $p(x_1)=q(x_1)$。同理对所有 $x_i$ 均有 $p(x_i)=q(x_i)$。一个 $n$ 次多项式在 $n+1$ 个点取值相同必恒等，矛盾。
\end{proof}

利用唯一性可得一些有用的推论。
\begin{theorem}
设 $f\in C_{2\pi}$ 为偶函数，$t\in T_n$ 是其最佳逼近，则 $t$ 也是偶函数。
\end{theorem}
\begin{proof}
由 $\|f(x)-t(x)\| \leqslant e_n(f)$ 得 $\|f(-x)-t(-x)\| \leqslant e_n(f)$，故 $\|f(x)-t(-x)\| \leqslant e_n(f)$，即 $s(x)=t(-x)$ 也是最佳逼近。由唯一性 $t(x)=s(x)=t(-x)$。
\end{proof}

切比雪夫定理完整刻画了最佳逼近的充要条件，但一般而言，直接求出给定函数的最佳逼近多项式并不容易。这促使人们寻找另一类方法：构造结构简单的线性算子序列，它们虽未必是“最佳”的，但能明确写出，且一致收敛于目标函数。 Korovkin 定理恰好为这类构造性算子的收敛性提供了统一判据。从下一章开始，我们将具体考察若干经典的线性正算子及其逼近性质。

\section{经典线性正算子与 Korovkin 定理}

Weierstrass 第一逼近定理和第二逼近定理分别肯定了闭区间上的连续函数可用多项式一致逼近，周期连续函数可用三角多项式一致逼近。然而，这两个定理本身并未提供统一的构造方法。本章将从具体的构造性算子入手，给出 Bernstein、Baskakov、Mirakyan 以及 Vallée-Poussin 算子的直接证明，随后引入 Korovkin 定理，将这些算子的收敛性纳入一个简洁的统一框架之中。

\subsection{Weierstrass 第一逼近定理与 Bernstein 多项式}

Weierstrass 第一逼近定理断言：闭区间上的任何连续函数都能用一列多项式一致逼近。用数学语言描述，就是
\[
\forall \epsilon > 0,\ \exists N,\ \forall n \geq N,\ \forall x \in [a, b],\ |f(x) - p_n(x)| < \epsilon,
\]
其中 $\{p_n(x)\}$ 是一列多项式。下面给出具体证明。

我们先考虑 $f(x)$ 在闭区间 $[0,1]$ 上的情形。使用 Bernstein 多项式来构造逼近函数：
\[
B_n^f(x) = \sum_{k=0}^{n} f\left(\frac{k}{n}\right) \binom{n}{k} x^k (1-x)^{n-k}.
\]
构造的动机与概率论中的二项分布密切相关：进行 $n$ 次独立试验，每次成功概率为 $p$，成功 $k$ 次的概率为
\[
P(X = k) = \binom{n}{k} p^k (1-p)^{n-k}.
\]
当 $k$ 接近 $np$ 时概率最大。因此 $B_n^f(x)$ 相当于对 $f$ 在 $k/n\approx x$ 处的取值进行加权平均，使得 $B_n^f(x)$ 能更好地逼近 $f(x)$。

为证明一致收敛，需要以下组合恒等式。

\begin{lemma}
\[
\sum_{k=0}^{n} \binom{n}{k} u^k v^{n-k} = (u + v)^n.
\]
\end{lemma}
\begin{lemma}
\[
\sum_{k=0}^{n} k \binom{n}{k} u^k v^{n-k} = nu(u + v)^{n-1}.
\]
\end{lemma}
\begin{lemma}
\[
\sum_{k=0}^{n} k^2 \binom{n}{k} u^k v^{n-k} = nu(u+v)^{n-1} + n(n-1)u^2(u+v)^{n-2}.
\]
\end{lemma}

这些可由二项式定理对 $u$ 求导得到。特别地，取 $u=v=1$ 可得推论：
\[
\sum_{k=0}^{n} k \binom{n}{k} = n \cdot 2^{n-1},\qquad
\sum_{k=0}^{n} k^2 \binom{n}{k} = n(n+1)2^{n-2}.
\]

代入 $u=x,\ v=1-x$，由引理1得
\[
\sum_{k=0}^{n} \binom{n}{k} x^k (1-x)^{n-k} = 1,
\]
并可计算出关键的方差型求和：
\begin{align*}
\sum_{k=0}^{n} (k-nx)^2 \binom{n}{k} x^k (1-x)^{n-k}
&= \sum_{k=0}^{n} k^2 \binom{n}{k} x^k (1-x)^{n-k} \\
&\quad - 2nx\sum_{k=0}^{n} k\binom{n}{k}x^k(1-x)^{n-k} + n^2x^2\sum_{k=0}^{n}\binom{n}{k}x^k(1-x)^{n-k} \\
&= [nx + n(n-1)x^2] - 2nx\cdot(nx) + n^2x^2 \\
&= nx(1-x).
\end{align*}

现在证明 $B_n^f(x)$ 一致收敛于 $f(x)$。

\begin{proof}
由于 $f$ 在 $[0,1]$ 上连续，故一致连续：
\[
\forall \epsilon>0,\ \exists\delta>0,\ \forall x,y\in[0,1],\ |x-y|<\delta \Rightarrow |f(x)-f(y)|<\epsilon.
\]
利用 $\sum_{k=0}^{n} \binom{n}{k} x^k (1-x)^{n-k} = 1$，可将 $f(x)$ 改写为
\[
f(x) = \sum_{k=0}^{n} f(x) \binom{n}{k} x^k (1-x)^{n-k},
\]
于是
\[
B_n^f(x)-f(x) = \sum_{k=0}^{n} \left(f\left(\frac{k}{n}\right)-f(x)\right) \binom{n}{k} x^k (1-x)^{n-k}.
\]
按 $|k/n-x|<\delta$ 和 $\ge\delta$ 拆分估计。

当 $|k/n-x|<\delta$ 时，$|f(k/n)-f(x)|<\epsilon$，故
\[
\left|\sum_{|k/n-x|<\delta} \cdots \right|
\le \sum_{|k/n-x|<\delta} \epsilon \binom{n}{k} x^k (1-x)^{n-k}
\le \epsilon \sum_{k=0}^{n} \binom{n}{k} x^k (1-x)^{n-k} = \epsilon.
\]

当 $|k/n-x|\ge\delta$ 时，设 $|f(x)|\le M$，则 $|f(k/n)-f(x)|\le 2M$。利用 $(k-nx)^2 \ge n^2\delta^2$ 放大：
\[
\left|\sum_{|k/n-x|\ge\delta} \cdots \right|
\le \frac{2M}{n^2\delta^2} \sum_{k=0}^{n} (k-nx)^2 \binom{n}{k} x^k (1-x)^{n-k}
= \frac{2M}{n^2\delta^2}\cdot nx(1-x).
\]
因为 $x(1-x)\le 1/4$，上式 $\le \frac{M}{2n\delta^2}$。

综合两部分得
\[
|B_n^f(x)-f(x)| < \epsilon + \frac{M}{2n\delta^2}.
\]
取 $N > \frac{M}{2\epsilon\delta^2}$，则当 $n>N$ 时 $|B_n^f(x)-f(x)| < 2\epsilon$，故 $B_n^f(x)$ 在 $[0,1]$ 上一致收敛于 $f(x)$。
\end{proof}

对于一般闭区间 $[a,b]$ 上的连续函数 $f(x)$，通过线性变换
\[
g(t) = f(a+(b-a)t),\quad t\in[0,1],
\]
将问题化归到 $[0,1]$ 上。$g(t)$ 仍连续，存在 Bernstein 多项式 $B_n^g(t)$ 一致收敛于 $g(t)$。将 $t=\frac{x-a}{b-a}$ 代回，得到
\[
B_n^f(x) = B_n^g\!\left( \frac{x-a}{b-a} \right)
= \sum_{k=0}^{n} f\!\left( a+\frac{k}{n}(b-a) \right)
\binom{n}{k} \left( \frac{x-a}{b-a} \right)^{\!k}
\left( 1-\frac{x-a}{b-a} \right)^{\!n-k},
\]
它在 $[a,b]$ 上一致收敛于 $f(x)$。这就完成了 Weierstrass 第一逼近定理的完整证明。


\subsection{Fejér 算子与 Weierstrass 第二逼近定理}

Weierstrass 第二逼近定理断言：任何以 $2\pi$ 为周期的连续函数都能用一列三角多项式一致逼近。
Fejér 给出了一个构造性的证明，其核心是对傅里叶级数的部分和取算术平均。

设 $f\in C_{2\pi}$，其傅里叶级数的前 $n+1$ 个部分和为 $S_0(f,x),S_1(f,x),\dots ,S_n(f,x)$。
定义
\[
\sigma_n(f,x)=\frac{S_0(f,x)+S_1(f,x)+\cdots +S_n(f,x)}{n+1},
\]
称为 \textbf{Fejér 算子}。每个 $\sigma_n(f,x)$ 都是次数不超过 $n$ 的三角多项式。

由傅里叶级数理论，部分和可表示为 Dirichlet 核的卷积：
\[
S_k(f,x)=\frac1\pi\int_{-\pi}^{\pi}f(x-t)D_k(t)\,dt,\qquad 
D_k(t)=\frac{\sin\bigl((k+\frac12)t\bigr)}{2\sin\frac{t}{2}}.
\]
代入平均得
\[
\sigma_n(f,x)=\frac1\pi\int_{-\pi}^{\pi}f(x-t)\mathcal{F}_n(t)\,dt,
\]
其中
\[
\mathcal{F}_n(t)=\frac1{n+1}\sum_{k=0}^{n}D_k(t)
=\frac{1}{2(n+1)}\frac{\sin^2\!\bigl(\frac{n+1}{2}t\bigr)}{\sin^2\frac{t}{2}}
\]
称为 \textbf{Fejér 核}。它有两个基本性质：
\[
\mathcal{F}_n(t)\ge 0,\qquad 
\frac1\pi\int_{-\pi}^{\pi}\mathcal{F}_n(t)\,dt=1.
\]
非负性来自表达式中分子分母均非负，归一性则因为每个 $D_k$ 的积分为 $\pi$。

下面证明 $\sigma_n(f,x)$ 一致收敛于 $f(x)$，证明结构与 Bernstein 多项式的情形完全平行。

\begin{proof}
由于 $f\in C_{2\pi}$，它在 $\mathbb{R}$ 上一致连续：
\[
\forall\varepsilon>0,\ \exists\delta>0,\ \forall x,y:\ |x-y|<\delta\Rightarrow |f(x)-f(y)|<\varepsilon.
\]
利用归一化条件，将 $f(x)$ 写成积分形式：
\[
f(x)=\frac1\pi\int_{-\pi}^{\pi}f(x)\mathcal{F}_n(t)\,dt,
\]
于是
\[
\sigma_n(f,x)-f(x)=\frac1\pi\int_{-\pi}^{\pi}\bigl[f(x-t)-f(x)\bigr]\mathcal{F}_n(t)\,dt.
\]
将积分区间按 $|t|<\delta$ 和 $\delta\le|t|\le\pi$ 拆成两部分。

\textbf{第一部分} $|t|<\delta$。此时 $|f(x-t)-f(x)|<\varepsilon$，结合 $\mathcal{F}_n(t)\ge0$，
\[
\frac1\pi\biggl|\int_{|t|<\delta}\bigl[f(x-t)-f(x)\bigr]\mathcal{F}_n(t)\,dt\biggr|
\le\frac\varepsilon\pi\int_{|t|<\delta}\mathcal{F}_n(t)\,dt
\le\frac\varepsilon\pi\int_{-\pi}^{\pi}\mathcal{F}_n(t)\,dt=\varepsilon.
\]

\textbf{第二部分} $\delta\le|t|\le\pi$。设 $M=\|f\|_\infty$，则 $|f(x-t)-f(x)|\le2M$。此时分母满足 $\sin^2\frac{t}{2}\ge\sin^2\frac\delta2$，故
\[
\mathcal{F}_n(t)=\frac{1}{2(n+1)}\frac{\sin^2\bigl(\frac{n+1}{2}t\bigr)}{\sin^2\frac{t}{2}}
\le\frac{1}{2(n+1)\sin^2\frac\delta2}.
\]
因此
\[
\frac1\pi\biggl|\int_{\delta\le|t|\le\pi}\bigl[f(x-t)-f(x)\bigr]\mathcal{F}_n(t)\,dt\biggr|
\le\frac{2M}{\pi}\int_{\delta\le|t|\le\pi}\mathcal{F}_n(t)\,dt
\le\frac{2M}{\pi}\cdot 2\pi\cdot\frac{1}{2(n+1)\sin^2\frac\delta2}
=\frac{2M}{(n+1)\sin^2\frac\delta2}.
\]

综合两部分，对任意 $x$ 一致地有
\[
|\sigma_n(f,x)-f(x)|
<\varepsilon+\frac{2M}{(n+1)\sin^2\frac\delta2}.
\]
取 $N$ 足够大，使得当 $n>N$ 时 $\frac{2M}{(n+1)\sin^2\frac\delta2}<\varepsilon$，则
\[
|\sigma_n(f,x)-f(x)|<2\varepsilon,
\]
故 $\sigma_n(f,x)$ 在 $[-\pi,\pi]$ 上一致收敛于 $f(x)$。由周期性，收敛在整个实轴上一致。
\end{proof}

以上证明充分体现了线性正算子与核的非负性在逼近论中的核心作用，与 Bernstein 多项式的证明如出一辙。


%\subsection{半无穷区间上的推广：Baskakov 与 Mirakyan 算子}

\subsubsection{Baskakov 线性正算子序列}

研究完闭区间 $[0,1]$ 上的 Weierstrass 第一逼近定理和 Bernstein 多项式后，一个自然的想法是：如果函数定义在半无穷区间 $[0, +\infty)$ 上，我们是否还能构造类似的函数列来一致逼近？这里单纯使用多项式是做不到的，需要引入一些更复杂的构造。

Baskakov 在 1957 年提出了如下线性正算子序列：
\[
\mathcal{B}_n(f;x) = \sum_{k=0}^{\infty} f\left( \frac{k}{n} \right) \binom{n+k-1}{k} x^k (1+x)^{-n-k},
\]
其中 $x \in [0, +\infty)$。这个构造的直观动机依然与概率论密切相关。在 Bernstein 多项式中，我们利用的是二项分布；而这里的核函数 $b_{n,k}(x) = \binom{n+k-1}{k} x^k (1+x)^{-n-k}$，本质上是取成功概率 $p = \frac{1}{1+x}$ 后的\textbf{负二项分布}。这意味着我们仍然在对函数 $f$ 进行某种“加权平均”，只不过求和是在无穷级数上进行的。

为了证明该算子一致收敛到 $f(x)$，我们需要建立类似于 Bernstein 多项式情形下的组合恒等式。首先，考虑广义二项式定理：当 $|z| < 1$ 时，
\[
(1-z)^{-n} = \sum_{k=0}^{\infty} \binom{n+k-1}{k} z^k.
\]
对这一展开式可作如下理解：
\[
(1-z)^{-n} = \left( \sum_{m=0}^{\infty} z^m \right)^n = \sum_{k=0}^{\infty} \binom{n+k-1}{k} z^k.
\]
基于此，我们通过对 $z$ 求导建立三个关键引理。

\begin{lemma}
\label{lem:baskakov1}
\[
\sum_{k=0}^{\infty} \binom{n+k-1}{k} x^k (1+x)^{-n-k} = 1.
\]
\end{lemma}
\begin{proof}
令 $z = \frac{x}{1+x}$（当 $x \geq 0$ 时 $0 \leq z < 1$），由广义二项式定理有 $\sum_{k=0}^{\infty} \binom{n+k-1}{k} z^k = (1-z)^{-n}$。两边同乘 $(1+x)^{-n}$ 并代回 $z$ 即得。该引理表明算子具有单位性。
\end{proof}

\begin{lemma}
\label{lem:baskakov2}
\[
\sum_{k=0}^{\infty} k \binom{n+k-1}{k} x^k (1+x)^{-n-k} = nx.
\]
\end{lemma}
\begin{proof}
对广义二项式定理两边关于 $z$ 求导，得 $n(1-z)^{-n-1} = \sum_{k=0}^{\infty} k \binom{n+k-1}{k} z^{k-1}$。两边同乘 $z(1+x)^{-n}$ 并代入 $z = \frac{x}{1+x}$ 即得。
\end{proof}

\begin{lemma}
\label{lem:baskakov3}
\[
\sum_{k=0}^{\infty} k^2 \binom{n+k-1}{k} x^k (1+x)^{-n-k} = n(n+1)x^2 + nx.
\]
\end{lemma}
\begin{proof}
对引理 \ref{lem:baskakov2} 对应的 $z$ 恒等式再次关于 $z$ 求导并同乘 $z$，或利用二阶导数性质即可导出。
\end{proof}

有了上述工具，我们可以仿照 Weierstrass 定理的证明思路，估计算子与 $f(x)$ 之间的偏差。

首先，利用引理 \ref{lem:baskakov1}--\ref{lem:baskakov3}，得到如下组合恒等式：
\begin{align*}
\sum_{k=0}^{\infty} (k-nx)^2 b_{n,k}(x)
&= \sum_{k=0}^{\infty} k^2 b_{n,k}(x) - 2nx \sum_{k=0}^{\infty} k b_{n,k}(x) + n^2x^2 \sum_{k=0}^{\infty} b_{n,k}(x) \\
&= [n(n+1)x^2 + nx] - 2nx(nx) + n^2x^2 \\
&= nx(1+x).
\end{align*}

下面考虑在任意给定的闭区间 $[0, a]$ 上的一致收敛性。设 $f(x)$ 在该区间上连续且有界（$|f(x)| \leq M$）。由一致连续性，对任给 $\epsilon > 0$，存在 $\delta > 0$，使得当 $|k/n - x| < \delta$ 时，有 $|f(k/n) - f(x)| < \epsilon$。

将求和按 $|k/n - x| < \delta$ 和 $|k/n - x| \geq \delta$ 分为两部分进行估计：
\begin{align*}
\left| \sum_{k=0}^{\infty} f\left(\frac{k}{n}\right) b_{n,k}(x) - f(x) \right|
&= \left| \sum_{k=0}^{\infty} [f(k/n) - f(x)] b_{n,k}(x) \right| \\
&\leq \sum_{|k/n - x| < \delta} \left|[f(k/n) - f(x)] b_{n,k}(x) \right| + \sum_{|k/n - x| \geq \delta} \left|[f(k/n) - f(x)] b_{n,k}(x) \right| \\
&< \epsilon \sum_{k=0}^{\infty} b_{n,k}(x) + \frac{2M}{n^2\delta^2} \sum_{k=0}^{\infty} (k-nx)^2 b_{n,k}(x) \\
&= \epsilon + \frac{2M \cdot nx(1+x)}{n^2\delta^2} = \epsilon + \frac{2Mx(1+x)}{n\delta^2}.
\end{align*}

在区间 $[0, a]$ 上，$x(1+x) \leq a(1+a)$ 有界。因此对于选定的 $\epsilon$ 和 $\delta$，只要 $n$ 足够大，使得 $n > \frac{2Ma(1+a)}{\epsilon\delta^2}$，偏差就能控制在 $2\epsilon$ 以内。这与 Bernstein 多项式的情形非常类似，本质上都是通过分段估计来控制误差。当然，也可以用 Korovkin 定理来处理这一问题。

综上所述，算子序列 $\mathcal{B}_n(f;x)$ 在 $[0, a]$ 上一致收敛到 $f(x)$。

\subsubsection{Mirakyan 算子序列}

设 $f$ 是 $\mathbb{R}$ 上的连续函数。定义 Mirakyan 算子序列
\[
M_n(f;x) = \sum_{k=0}^{\infty} f\left(\frac{k}{n}\right) \frac{n^k}{k!} x^k e^{-nx},
\]
下面证明 $M_n(f;x)$ 在任意闭区间 $[0, a]$ 上一致收敛于 $f$。

核函数 $\dfrac{(nx)^k}{k!}e^{-nx}$ 即为 Poisson 分布取参数 $\lambda = nx$ 的概率质量函数。Poisson 分布与二项分布关系密切：在 Weierstrass 第一逼近定理中，我们通过二项分布构造了 Bernstein 多项式 $B_n^f(x) = \sum_{k=0}^{n} f(\frac{k}{n}) \binom{n}{k} x^k (1-x)^{n-k}$，而这里相当于将核函数替换为 Poisson 分布。注意到 Poisson 极限公式
\[
\lim_{m \to \infty} \binom{m}{k} \left(\frac{nx}{m}\right)^k \left(1 - \frac{nx}{m}\right)^{m-k} = \frac{(nx)^k}{k!} e^{-nx},
\]
因此 Mirakyan 算子的收敛性也可以用来证明 Weierstrass 第一逼近定理。

由 $f$ 在 $[0,a]$ 上的一致连续性，$\forall \epsilon >0,\exists \delta>0$，使得当 $|x-y|<\delta$ 且 $x,y \in [0,a]$ 时，$|f(x)-f(y)|<\epsilon$。与 Weierstrass 定理的证明类似，这里同样有一个基本恒等式：
\[
\sum_{k=0}^{\infty} \frac{n^k}{k!} x^k e^{-nx} = 1.
\]
这由指数函数的 Taylor 展开可直接验证：
\[
\sum_{k=0}^{\infty} \frac{n^k}{k!} x^k e^{-nx}
= e^{-nx} \sum_{k=0}^{\infty} \frac{(nx)^k}{k!}
= e^{-nx} \cdot e^{nx} = 1.
\]

令 $|M_n(f;x)-f(x)| = \left|\sum_{k=0}^{\infty} \left(f\left(\frac{k}{n}\right) - f(x)\right)\frac{n^k}{k!} x^k e^{-nx}\right|$，则
\begin{align*}
|M_n(f;x)-f(x)|
&= \left| \left(\sum_{|k/n - x| < \delta} + \sum_{|k/n - x| \geq \delta}\right) \left(f\left(\frac{k}{n}\right) - f(x)\right) \frac{n^k}{k!} x^k e^{-nx} \right|.
\end{align*}

对于第一部分，利用一致连续性，
\begin{align*}
\left| \sum_{|k/n - x| < \delta} \left(f\left(\frac{k}{n}\right) - f(x)\right) \frac{n^k}{k!} x^k e^{-nx} \right|
&\leq \sum_{|k/n - x| < \delta} \left| f\left(\frac{k}{n}\right) - f(x) \right| \frac{n^k}{k!} x^k e^{-nx} \\
&\leq \sum_{k=0}^{\infty} \epsilon \frac{n^k}{k!} x^k e^{-nx}
= \epsilon \sum_{k=0}^{\infty} \frac{n^k}{k!} x^k e^{-nx} = \epsilon.
\end{align*}

对于第二部分，由于 $f$ 在 $[0,a]$ 上连续故有界，设 $\|f\| \leq M$，则 $|f(\frac{k}{n}) - f(x)| \leq 2M$。于是
\begin{align*}
\left| \sum_{|k/n - x| \geq \delta} \left(f\left(\frac{k}{n}\right) - f(x)\right) \frac{n^k}{k!} x^k e^{-nx} \right|
&\leq \sum_{|k/n - x| \geq \delta} 2M \frac{n^k}{k!} x^k e^{-nx}.
\end{align*}

由条件 $|\frac{k-nx}{n}| \geq \delta$ 可得 $(k-nx)^2 \geq (n\delta)^2$，因此
\begin{align*}
\sum_{|k/n - x| \geq \delta} 2M \frac{n^k}{k!} x^k e^{-nx}
&\leq \sum_{k=0}^{\infty} \frac{(k-nx)^2}{(n\delta)^2} 2M \frac{n^k}{k!} x^k e^{-nx} \\
&= \frac{2M}{n^2\delta^2} \sum_{k=0}^{\infty} (k-nx)^2 \frac{n^k}{k!} x^k e^{-nx}.
\end{align*}

为计算上述求和，利用指数函数幂级数的性质：
\[
\sum_{k=0}^{\infty} \frac{x^k}{k!} = e^x,
\]
\[
\sum_{k=0}^{\infty} k \cdot \frac{x^k}{k!}
= \sum_{k=1}^{\infty} \frac{x^k}{(k-1)!}
= \sum_{l=0}^{\infty} \frac{x^{l+1}}{l!}
= x e^{x},
\]
\[
\sum_{k=0}^{\infty} k^2 \frac{x^k}{k!}
= \sum_{k=1}^{\infty} k \cdot \frac{x^k}{(k-1)!}
= \sum_{k=1}^{\infty} (k-1) \cdot \frac{x^k}{(k-1)!} + \sum_{k=1}^{\infty} \frac{x^k}{(k-1)!}
= x^2 e^x + x e^x.
\]

代入可得
\begin{align*}
\frac{2M}{n^2\delta^2} \sum_{k=0}^{\infty} (k-nx)^2 \frac{n^k}{k!} x^k e^{-nx}
&= \frac{2M}{n^2\delta^2} \sum_{k=0}^{\infty} (k^2 - 2nxk + n^2x^2) \frac{n^k}{k!} x^k e^{-nx} \\
&= \frac{2M}{n^2\delta^2} \left( (nx)^2 e^{nx} + (nx) e^{nx} - 2nx(nx) e^{nx} + (nx)^2 e^{nx} \right) e^{-nx} \\
&= \frac{2Mx}{n\delta^2}.
\end{align*}

取 $N = \left[\frac{2Ma}{\epsilon \delta^2}\right]+1$，则当 $n > N$ 时，上式小于 $\epsilon$。综合两部分估计，得到
\[
|M_n(f;x)-f(x)| < 2\epsilon
\]
对任意 $x \in [0,a]$ 成立。因此 $M_n(f;x)$ 在 $[0,a]$ 上一致收敛于 $f(x)$。


%\input{chapter3_korovkin/theorem/vallee_poussin.tex}

\subsection{Korovkin 定理——线性正算子收敛的统一判据}

前面几节我们逐一给出了 Bernstein、Baskakov、Mirakyan 以及 V-P 算子的构造与收敛性证明。回顾这些证明，不难发现它们共享着几乎完全相同的论证结构：利用核函数的非负性与归一性，将逼近误差按远近拆分为两部分，近处用一致连续性控制，远处用函数的界与核函数的方差型求和控制，最终令 $n\to\infty$ 得到一致收敛。

这种高度雷同的证明模式暗示着背后存在一个统一的原理。Korovkin 在 20 世纪 50 年代发现了这一原理：对于线性正算子，其一致收敛性完全由三个“测试函数”的逼近行为所决定。在代数情形下，测试函数为 $1$, $x$, $x^2$；在三角情形下，测试函数为 $1$, $\cos x$, $\sin x$。本节将阐述并证明这两个 Korovkin 定理，并展示如何用它们简洁地统一处理前述所有算子。

\subsubsection{Korovkin 定理（代数多项式情形）}

\begin{theorem}[Korovkin 定理，代数多项式情形]
设线性正算子 $L_n: C[a,b]\to C[a,b]$ 满足
\[
L_n(1;x)=1+\alpha_n(x),\quad L_n(t;x)=x+\beta_n(x),\quad L_n(t^2;x)=x^2+\theta_n(x),
\]
其中 $\alpha_n,\beta_n,\theta_n$ 在 $[a,b]$ 上一致收敛于 $0$。则对任意 $f\in C[a,b]$，$L_n(f;x)$ 一致收敛于 $f(x)$。
\end{theorem}

\begin{proof}
由 $f$ 在 $[a,b]$ 上一致连续且 $\|f\|\le M$，$\forall \epsilon>0,\exists\delta>0$，当 $|x-y|<\delta$ 时 $|f(x)-f(y)|<\epsilon$。对任意 $x,t\in[a,b]$ 有
\[
-\frac{2M}{\delta^2}(t-x)^2 - \epsilon \le f(t)-f(x) \le \frac{2M}{\delta^2}(t-x)^2 + \epsilon.
\]
作用线性正算子 $L_n$（关于变量 $t$，$x$ 视为参数），利用线性性与非负性保持不等式方向，得
\[
-\frac{2M}{\delta^2}L_n((t-x)^2;x)-\epsilon L_n(1;x)
\le L_n(f;x)-f(x)L_n(1;x)
\le \frac{2M}{\delta^2}L_n((t-x)^2;x)+\epsilon L_n(1;x).
\]
计算
\begin{align*}
L_n((t-x)^2;x)
&= L_n(t^2;x) - 2xL_n(t;x) + x^2L_n(1;x)\\
&= (x^2+\theta_n(x)) - 2x(x+\beta_n(x)) + x^2(1+\alpha_n(x))\\
&= \theta_n(x) - 2x\beta_n(x) + x^2\alpha_n(x),
\end{align*}
它一致收敛于 $0$。又 $L_n(1;x)=1+\alpha_n(x)\to 1$。所以
\[
L_n(f;x) - f(x) = \big(L_n(f;x)-f(x)L_n(1;x)\big) + f(x)\alpha_n(x) \to 0
\]
一致成立。
\end{proof}

\subsubsection{Korovkin 定理（三角多项式情形）}

对于周期函数，有完全平行的结论。

\begin{theorem}[Korovkin 定理，三角多项式情形]
设线性正算子 $L_n: C_{2\pi}\to C_{2\pi}$ 满足
\[
L_n(1;x)=1+\alpha_n(x),\quad L_n(\cos t;x)=\cos x+\beta_n(x),\quad L_n(\sin t;x)=\sin x+\theta_n(x),
\]
且 $\alpha_n,\beta_n,\theta_n$ 一致收敛于 $0$，则对任意 $f\in C_{2\pi}$，$L_n(f;x)$ 一致收敛于 $f(x)$。
\end{theorem}

\begin{proof}
证明思路与代数情形类似。由 $f$ 在 $\R$ 上一致连续且以 $2\pi$ 为周期，$\forall\epsilon>0,\exists\delta>0$，当 $|x-y|<\delta$ 时 $|f(x)-f(y)|<\epsilon$。对任意 $x,t\in\R$，可用三角不等式得到
\[
-\frac{2M}{\sin^2(\delta/2)}\sin^2\frac{t-x}{2} - \epsilon \le f(t)-f(x) \le \frac{2M}{\sin^2(\delta/2)}\sin^2\frac{t-x}{2} + \epsilon.
\]
利用恒等式
\[
\sin^2\frac{t-x}{2} = \frac{1}{2}\big(1 - \cos t\cos x - \sin t\sin x\big),
\]
对 $L_n$ 作用后，$L_n(\sin^2\frac{t-x}{2};x)$ 可展开为 $L_n(1;x)$, $L_n(\cos t;x)$, $L_n(\sin t;x)$ 的线性组合。由定理条件，这三项分别一致收敛于 $1$, $\cos x$, $\sin x$，因此 $L_n(\sin^2\frac{t-x}{2};x)$ 一致收敛于 $0$。其余步骤与代数情形完全相同。
\end{proof}

\subsubsection{统一验证：回到经典算子}

Korovkin 定理的美妙之处在于，它将原本需要逐一进行的分段估计，统一简化为对三个测试函数的验证。下面以 Bernstein 多项式和 Fejér 算子为例，展示这种统一性在代数情形与三角情形中的平行应用。

\textbf{Bernstein 多项式}：对 $f\in C[0,1]$，$B_n(f;x)=\sum_{k=0}^{n} f(\frac{k}{n})\binom{n}{k}x^{k}(1-x)^{n-k}$ 是线性正算子。直接计算三个测试函数：
\[
B_n(1;x)=1,\qquad B_n(t;x)=x,\qquad B_n(t^{2};x)=x^{2}+\frac{x(1-x)}{n}.
\]
因此 $\alpha_n(x)=0$，$\beta_n(x)=0$，$\theta_n(x)=\frac{x(1-x)}{n}\rightrightarrows 0$。由 Korovkin 定理即得 $B_n(f;x)$ 一致收敛于 $f(x)$。

\textbf{Fejér 算子}：对周期函数 $f\in C_{2\pi}$，Fejér 算子
\[
\sigma_n(f,x)=\frac{1}{\pi}\int_{-\pi}^{\pi}f(x-t)\mathcal{F}_n(t)\,dt,\qquad
\mathcal{F}_n(t)=\frac{1}{2(n+1)}\frac{\sin^{2}\bigl(\frac{n+1}{2}t\bigr)}{\sin^{2}\frac{t}{2}}
\]
是线性正算子（积分核非负）。利用三角形式的 Korovkin 定理，只需验证测试函数 $1,\cos t,\sin t$。

由归一化条件即得 $\sigma_n(1,x)=1$。对于 $\cos t$，利用 $\cos(x-t)=\cos x\cos t+\sin x\sin t$ 以及 $\mathcal{F}_n$ 为偶函数，有
\[
\sigma_n(\cos t,x)=\frac{1}{\pi}\int_{-\pi}^{\pi}\cos(x-t)\mathcal{F}_n(t)\,dt
=\cos x\cdot\frac{1}{\pi}\int_{-\pi}^{\pi}\cos t\,\mathcal{F}_n(t)\,dt.
\]
由 Dirichlet 核的性质，$\frac{1}{\pi}\int_{-\pi}^{\pi}\cos t\,D_k(t)\,dt=1$（$k\ge 1$）而 $k=0$ 时为 $0$，因此
\[
\frac{1}{\pi}\int_{-\pi}^{\pi}\cos t\,\mathcal{F}_n(t)\,dt
=\frac{1}{n+1}\sum_{k=0}^{n}\frac{1}{\pi}\int_{-\pi}^{\pi}\cos t\,D_k(t)\,dt
=\frac{n}{n+1}=1-\frac{1}{n+1}.
\]
于是 $\sigma_n(\cos t,x)=\frac{n}{n+1}\cos x=\cos x-\frac{1}{n+1}\cos x$，余项一致趋于 $0$。同理，由奇偶性可得
\[
\sigma_n(\sin t,x)=\sin x\cdot\frac{1}{\pi}\int_{-\pi}^{\pi}\cos t\,\mathcal{F}_n(t)\,dt
=\frac{n}{n+1}\sin x=\sin x-\frac{1}{n+1}\sin x,
\]
余项亦一致趋于 $0$。由三角形式 Korovkin 定理立即得到 $\sigma_n(f,x)\rightrightarrows f(x)$，这样 Weierstrass 第二逼近定理作为 Korovkin 定理的直接推论而成立，与 Bernstein 多项式的情形形成完美对照。

\subsection{小结}

本章从具体的构造性算子出发，给出了 Weierstrass 两个逼近定理的完整证明，随后引入 Korovkin 定理，揭示了这些算子收敛性的共同本质：只需验证三个测试函数的矩条件，即可保证对一切连续函数的一致收敛。Bernstein 多项式的二项分布核、Baskakov 算子的负二项分布核、Mirakyan 算子的 Poisson 分布核、Vallée-Poussin 算子的幂次余弦核，虽然形式各异，但在 Korovkin 的框架下，它们的收敛性验证都归结为完全相同的简单计算。

然而，Korovkin 定理只回答了“是否收敛”的问题，并未涉及逼近的速率。要刻画“收敛有多快”，需要将线性算子的误差与最佳逼近误差联系起来，这正是下一章 Lebesgue 不等式的任务。

\section{Lebesgue不等式与三角多项式算子的Lebesgue常数}

在上一章中，我们借助Korovkin定理建立了线性正算子的一致收敛性。但收敛性本身只告诉我们“当$n$足够大时误差会趋于零”，并没有给出任何关于“收敛有多快”的信息。要讨论逼近速率，就必须引入一个连接“构造性算子”与“最佳逼近”之间的桥梁——Lebesgue不等式。

\subsection{Lebesgue不等式的意义与证明}

设$\mathcal{T}_n$为次数不超过$n$的三角多项式全体。对于$f\in C_{2\pi}$，记其最佳逼近误差为
\[
E_n(f)=\inf_{T\in\mathcal{T}_n}\|f-T\|_\infty.
\]
任给一个线性算子$L_n:C_{2\pi}\to\mathcal{T}_n$，我们希望用$E_n(f)$来控制实际逼近误差$\|f-L_nf\|_\infty$。如果$L_n$恰好是到$\mathcal{T}_n$上的投影（即对任何$T\in\mathcal{T}_n$有$L_nT=T$），那么下面的简单不等式就给出了答案。

\begin{theorem}[Lebesgue不等式]
设$L_n:C_{2\pi}\to\mathcal{T}_n$是一个线性算子，满足对任意$T\in\mathcal{T}_n$有$L_nT=T$。令
\[
\Lambda_n=\|L_n\|=\sup_{\|f\|_\infty=1}\|L_nf\|_\infty
\]
为算子的范数，称为\textbf{Lebesgue常数}。则对任何$f\in C_{2\pi}$，有
\[
\|f-L_nf\|_\infty\le (1+\Lambda_n)E_n(f).
\]
\end{theorem}

\begin{proof}
由最佳逼近的定义，存在$T^*\in\mathcal{T}_n$使得$\|f-T^*\|_\infty=E_n(f)$。利用$L_nT^*=T^*$以及$L_n$的线性性，
\[
f-L_nf=(f-T^*)-L_n(f-T^*).
\]
取范数，由三角不等式和算子范数的定义即得
\[
\|f-L_nf\|_\infty\le\|f-T^*\|_\infty+\|L_n(f-T^*)\|_\infty
\le E_n(f)+\Lambda_n\|f-T^*\|_\infty=(1+\Lambda_n)E_n(f).
\]
\end{proof}

这个不等式的重要性在于：它把用线性算子得到的实际逼近误差，分离成了两个因子的乘积——$(1+\Lambda_n)$是算子本身带来的“放大”，而$E_n(f)$则是该函数类的内在最佳逼近能力。$\Lambda_n$越小（越接近$1$），算子就越“不损失”最佳逼近的精度；反之若$\Lambda_n$随$n$增长而急剧增大，即使$E_n(f)$很小，实际误差也可能被大幅放大。

接下来我们考察三角逼近中两个最重要的算子——傅里叶部分和算子与Fejér算子，精确计算它们的Lebesgue常数。一个发散，一个恒为$1$，对比极富启发性。

\subsection{傅里叶部分和算子的Lebesgue常数}

对$f\in C_{2\pi}$，其傅里叶级数的$n$阶部分和为
\[
S_n(f,x)=\frac{a_0}{2}+\sum_{k=1}^{n}\bigl(a_k\cos kx+b_k\sin kx\bigr),
\]
其中$a_k,b_k$为傅里叶系数。算子$S_n$显然是线性的，并且次数不超过$n$的三角多项式在$S_n$作用下不变（因为它的傅里叶级数就是它自身）。因此$S_n$满足Lebesgue不等式的条件。

利用Dirichlet核
\[
D_n(t)=\frac12+\sum_{k=1}^{n}\cos kt=\frac{\sin\bigl((n+\frac12)t\bigr)}{2\sin\frac{t}{2}},
\]
部分和可写成卷积形式
\[
S_n(f,x)=\frac1\pi\int_{-\pi}^{\pi}f(x-t)D_n(t)\,dt.
\]
于是算子范数等于Dirichlet核的$L_1$范数：
\[
\Lambda_n=\|S_n\|=\frac1\pi\int_{-\pi}^{\pi}|D_n(t)|\,dt.
\]

\begin{theorem}[傅里叶部分和的Lebesgue常数的渐近阶]
存在常数$C>0$，使得
\[
\Lambda_n=\frac{4}{\pi^2}\ln n+O(1),\qquad n\to\infty.
\]
\end{theorem}

\begin{proof}
由$D_n$的偶性，只需考虑$[0,\pi]$上的积分：
\[
\Lambda_n=\frac2\pi\int_0^\pi|D_n(t)|\,dt
=\frac1\pi\int_0^\pi\frac{|\sin((n+\frac12)t)|}{\sin\frac{t}{2}}\,dt.
\]
令$m=n+\frac12$。在$[0,\pi]$上，利用$\sin\frac{t}{2}\ge\frac{t}{\pi}$（由正弦的凹性）可得上界
\[
\Lambda_n\le\frac1\pi\int_0^\pi\frac{|\sin mt|}{t/\pi}\,dt
=\int_0^\pi\frac{|\sin mt|}{t}\,dt.
\]
另一方面，由$\sin\frac{t}{2}\le\frac{t}{2}$得下界
\[
\Lambda_n\ge\frac1\pi\int_0^\pi\frac{|\sin mt|}{t/2}\,dt
=\frac2\pi\int_0^\pi\frac{|\sin mt|}{t}\,dt.
\]
因此只要处理积分
\[
I_m:=\int_0^\pi\frac{|\sin mt|}{t}\,dt
=\int_0^{m\pi}\frac{|\sin u|}{u}\,du\qquad(u=mt).
\]
将$[0,m\pi]$按$\sin u$的零点分为$m$个长度为$\pi$的小区间：
\[
I_m=\sum_{k=0}^{m-1}\int_{k\pi}^{(k+1)\pi}\frac{|\sin u|}{u}\,du.
\]
对于$k\ge1$，当$u\in[k\pi,(k+1)\pi]$时，$\frac1{(k+1)\pi}\le\frac1u\le\frac1{k\pi}$，而$\int_{k\pi}^{(k+1)\pi}|\sin u|\,du=2$。因此
\[
\frac{2}{(k+1)\pi}\le\int_{k\pi}^{(k+1)\pi}\frac{|\sin u|}{u}\,du\le\frac{2}{k\pi}.
\]
$k=0$的项$\int_0^\pi\frac{\sin u}{u}\,du$为常数，记作$A$。于是
\[
I_m=A+\frac{2}{\pi}\sum_{k=1}^{m-1}\theta_k,\qquad \frac{1}{k+1}\le\theta_k\le\frac1k.
\]
而$\sum_{k=1}^{m-1}\frac1k=\ln m+O(1)$，且$\sum_{k=1}^{m-1}(\frac1k-\frac1{k+1})=1-\frac1m$为有界量，故
\[
\sum_{k=1}^{m-1}\theta_k=\ln m+O(1).
\]
因此$I_m=\frac{2}{\pi}\ln m+O(1)$。代入$\Lambda_n$的上、下界表达式，注意$m\sim n$，便得到
\[
\Lambda_n=\frac{4}{\pi^2}\ln n+O(1).
\]
\end{proof}

对数发散虽然缓慢，但毕竟趋于无穷。由Banach–Steinhaus定理，这意味着存在连续函数其傅里叶级数不一致收敛（甚至在某些点发散）。这正是$\Lambda_n$无界所带来的本质障碍。

\subsection{Fejér算子的Lebesgue常数}

Fejér算子定义为部分和的算术平均：
\[
\sigma_n(f,x)=\frac{S_0(f,x)+S_1(f,x)+\cdots+S_n(f,x)}{n+1}.
\]
它也可以写成卷积形式
\[
\sigma_n(f,x)=\frac1\pi\int_{-\pi}^{\pi}f(x-t)\mathcal{F}_n(t)\,dt,
\]
其中Fejér核为
\[
\mathcal{F}_n(t)=\frac1{n+1}\sum_{k=0}^{n}D_k(t)
=\frac{1}{2(n+1)}\frac{\sin^2\!\bigl(\frac{n+1}{2}t\bigr)}{\sin^2\frac{t}{2}}.
\]
注意$\mathcal{F}_n(t)\ge0$且$\frac1\pi\int_{-\pi}^{\pi}\mathcal{F}_n(t)dt=1$。

与$S_n$不同，$\sigma_n$并不是到$\mathcal{T}_n$的投影（一个$T\in\mathcal{T}_n$的Fejér平均并不一定等于$T$自身），因此不能直接套用定理形式中的Lebesgue不等式。但我们仍然可以计算$\sigma_n$作为线性算子的范数，这个范数同样扮演着刻画稳定性的角色。

\begin{theorem}[Fejér算子的范数]
对任意$n$，有$\|\sigma_n\|=1$。
\end{theorem}

\begin{proof}
对任意$f$，由于$\mathcal{F}_n\ge0$，
\[
|\sigma_n(f,x)|\le\frac1\pi\int_{-\pi}^{\pi}|f(x-t)|\mathcal{F}_n(t)\,dt
\le\|f\|_\infty\cdot\frac1\pi\int_{-\pi}^{\pi}\mathcal{F}_n(t)\,dt=\|f\|_\infty.
\]
故$\|\sigma_n\|\le1$。另一方面，取$f\equiv1$，则$\sigma_n(1,x)\equiv1$，达到等号，因此$\|\sigma_n\|=1$。
\end{proof}

这一结果与傅里叶部分和形成鲜明对照：Fejér算子的范数不仅不随$n$增大，而且恒为最小值$1$。在下一章中，我们将利用连续模和Jackson定理，给出这两种算子逼近阶的精确估计，并揭示范数差异带来的本质影响。



\section{逼近阶的估计：Jackson定理与Zygmund类的应用}

有了Korovkin定理保证的收敛性和Lebesgue常数提供的稳定性度量，我们现在进入定量逼近的核心——逼近阶的计算。这里的思路是：先用Jackson定理把最佳逼近误差$E_n(f)$与函数的光滑性（连续模）挂钩，再结合上一章的范数信息，对傅里叶部分和与Fejér算子分别得出具体收敛速率。最后，引入Zygmund类作为统一描述光滑性的语言，给出逼近阶的充要条件，并重新审视$|x|$和$|\sin x|$这两个经典例子。

\subsection{连续模与Jackson定理}

对于$f\in C_{2\pi}$，其连续模定义为
\[
\omega(f,\delta)=\sup_{|h|\le\delta}\max_{x}|f(x+h)-f(x)|,\qquad\delta>0.
\]
连续模刻画了函数的一致连续性：$\omega(f,\delta)\to0$当$\delta\to0$，且若$f\in\mathrm{Lip}\,\alpha$（$0<\alpha\le1$），则$\omega(f,\delta)\le M\delta^\alpha$。

逼近论中一个里程碑式的结果是Jackson定理，它用连续模给出了最佳逼近误差的精确上界。

\begin{theorem}[Jackson定理，周期形式]
存在绝对常数$C>0$，使得对任意$f\in C_{2\pi}$和任意正整数$n$，有
\[
E_n(f)\le C\,\omega\!\left(f,\frac1n\right).
\]
\end{theorem}

这个定理的证明通常借助Jackson核（或Fejér核的适当幂次）构造出具体的逼近多项式，细节可参见标准教材。这里我们直接使用这个结果。Jackson定理表明，函数越光滑（连续模越小），其最佳逼近误差衰减得越快。特别地，若$f\in\mathrm{Lip}\,\alpha$，则
\[
E_n(f)=O(n^{-\alpha}).
\]

\subsection{傅里叶部分和的逼近阶}

对傅里叶部分和$S_n(f)$，直接应用Lebesgue不等式和Jackson定理，便得到
\[
\|f-S_n(f)\|_\infty\le (1+\Lambda_n)E_n(f)
\le C(1+\Lambda_n)\,\omega\!\left(f,\frac1n\right).
\]
由于$\Lambda_n\sim\frac{4}{\pi^2}\ln n$，这意味着即使$f$非常光滑，前面的对数因子也会让误差上界比最佳可能多一个对数增长。更精确地说，对$f\in\mathrm{Lip}\,\alpha$，有
\[
\|f-S_n(f)\|_\infty=O\!\left(\frac{\ln n}{n^\alpha}\right).
\]
对数项的存在是无法消除的：确实存在$\mathrm{Lip}\,\alpha$函数，其傅里叶部分和的误差能达到$\asymp n^{-\alpha}\ln n$。因此$S_n$不是最优的线性逼近算子，其“额外对数代价”正是无界Lebesgue常数的直接后果。

\subsection{Fejér算子的逼近阶}

Fejér算子的情况要稍微妙一些。由于它不具备投影性质，不能直接套用$(1+\Lambda_n)E_n(f)$的形式。但我们可以直接用算子本身的结构来估计它与$f$的差距，并借助连续模得出Jackson型的界。

对$f\in C_{2\pi}$，将$\sigma_n(f,x)-f(x)$写成
\[
\sigma_n(f,x)-f(x)=\frac1\pi\int_{-\pi}^{\pi}[f(x-t)-f(x)]\mathcal{F}_n(t)\,dt.
\]
由于$\mathcal{F}_n\ge0$，模仿Korovkin定量估计的标准手法，对任意$\delta>0$，
\[
\begin{aligned}
|\sigma_n(f,x)-f(x)|
&\le\frac1\pi\int_{|t|<\delta}\omega(f,\delta)\mathcal{F}_n(t)\,dt
+\frac1\pi\int_{\delta\le|t|\le\pi}2\|f\|_\infty\mathcal{F}_n(t)\,dt\\
&\le\omega(f,\delta)+\frac{2\|f\|_\infty}{\pi}\int_{\delta\le|t|\le\pi}\mathcal{F}_n(t)\,dt.
\end{aligned}
\]
利用$\mathcal{F}_n(t)\le\frac{1}{2(n+1)\sin^2(\delta/2)}$（对$|t|\ge\delta$）以及区间长度不超过$2\pi$，可得
\[
\frac1\pi\int_{\delta\le|t|\le\pi}\mathcal{F}_n(t)\,dt
\le\frac{2}{(n+1)\sin^2(\delta/2)}.
\]
取$\delta=n^{-1/2}$，则$\sin(\delta/2)\sim\delta/2$，于是上述尾部不超过$O(1/n\cdot 1/\delta^2)=O(1)$？仔细核算：当$\delta=n^{-1/2}$时，$\sin^2(\delta/2)\approx \delta^2/4=n^{-1}/4$，因此
\[
\frac{1}{(n+1)\sin^2(\delta/2)}\approx\frac{4}{n+1}\cdot n=O(1).
\]
这无法让尾部随$n$增大而趋于零。这说明取$\delta$为$n^{-1/2}$还不够小。更精细的估计需要利用Fejér核在无穷远处的衰减，或者直接采用经典的Korovkin定理的“测试函数”技巧。

事实上，上一章Korovkin定理的证明过程中，我们曾提取出一个一般性的定量不等式：对线性正算子$L_n$，
\[
\|L_nf-f\|_\infty\le 2\omega\!\left(f,\sqrt{\|\delta_n\|_\infty}\right),
\]
其中$\delta_n(x)=L_n((t-x)^2;x)$。对Fejér算子，可以用三角形式类似推导。考虑周期情形的“测试函数”$1,\cos t,\sin t$，其对应的中心矩为
\[
\sigma_n(\sin^2((t-x)/2);x)=\frac12[1-\sigma_n(\cos(t-x);x)].
\]
通过计算傅里叶系数可得$\sigma_n(\cos(t-x);x)=\cos x\cdot\frac1\pi\int_{-\pi}^{\pi}\cos t\,\mathcal{F}_n(t)dt+\sin x\cdot\frac1\pi\int_{-\pi}^{\pi}\sin t\,\mathcal{F}_n(t)dt$。由于$\mathcal{F}_n$是偶函数，正弦项积分为零；余弦项积分可由Dirichlet核积分得到：
\[
\frac1\pi\int_{-\pi}^{\pi}\cos t\,\mathcal{F}_n(t)dt
=\frac1{n+1}\sum_{k=0}^{n}\frac1\pi\int_{-\pi}^{\pi}\cos t\,D_k(t)dt
=\frac1{n+1}\sum_{k=0}^{n}\left(1-\frac{1}{k+1}\right)=1-\frac{H_{n+1}}{n+1},
\]
这里$H_{n+1}$是调和数。于是$\sigma_n(\cos(t-x);x)\approx\left(1-\frac{\ln n}{n}\right)\cos x$，进而
\[
\delta_n(x)=\sigma_n(2\sin^2((t-x)/2);x)=1-\sigma_n(\cos(t-x);x)\approx\frac{\ln n}{n}\cos x.
\]
这导致$\sqrt{\|\delta_n\|_\infty}\sim\sqrt{\frac{\ln n}{n}}$，从而给出估计
\[
\|\sigma_nf-f\|_\infty\le 2\omega\!\left(f,\sqrt{\frac{\ln n}{n}}\right).
\]
对$\mathrm{Lip}\,\alpha$函数，得到$\|\sigma_nf-f\|_\infty=O((\ln n)^{\alpha/2}n^{-\alpha/2})$。这一收敛速率比Bernstein算子的$O(n^{-\alpha/2})$多了一个对数因子，但本质上仍然是$n^{-\alpha/2}$量级。若直接使用Fejér算子的具体卷积结构，改进上述分析可得最优估计
\[
\|\sigma_nf-f\|_\infty\le C\,\omega\!\left(f,\frac1{\sqrt{n}}\right),
\]
这就是Fejér算子的Jackson型不等式。下面我们给出一个更为直接且避免对数因子的证明框架。

\begin{proof}[Fejér算子Jackson型不等式的证明概略]
利用连续模的性质，对任意$\delta>0$，
\[
\begin{aligned}
|\sigma_n(f,x)-f(x)|
&\le\frac1\pi\int_{-\pi}^{\pi}\omega(f,|t|)\mathcal{F}_n(t)\,dt\\
&\le\omega(f,\delta)+\frac{2\|f\|_\infty}{\pi}\int_{\delta\le|t|\le\pi}\mathcal{F}_n(t)\,dt.
\end{aligned}
\]
现在关键是要更精确地控制尾部积分。利用$\mathcal{F}_n(t)\le\frac{\pi^2}{2(n+1)t^2}$（由$\sin u\ge\frac{2}{\pi}u$在$[0,\pi/2]$上的估计可得），则当$\delta\le\pi$时，
\[
\int_{\delta\le|t|\le\pi}\mathcal{F}_n(t)\,dt\le 2\int_\delta^\pi\frac{\pi^2}{2(n+1)t^2}dt\le\frac{\pi^2}{n+1}\int_\delta^\infty\frac{dt}{t^2}=\frac{\pi^2}{(n+1)\delta}.
\]
选择$\delta=n^{-1/2}$，则尾部为$O(n^{-1/2})$，这已经能够与$\omega(f,\delta)$匹配。但更精细的平衡可取$\delta$使得$\omega(f,\delta)$与尾部同阶。若$\omega(f,\delta)\le M\delta$，取$\delta=n^{-1/2}$可得$O(n^{-1/2})$；对于一般$\omega(f,\delta)$，取$\delta=n^{-1/2}$，则$\omega(f,n^{-1/2})$控制主项，尾部为$O(n^{-1/2})$，因此最终有
\[
\|\sigma_nf-f\|_\infty\le C\omega\!\left(f,\frac1{\sqrt{n}}\right).
\]
\end{proof}

因此，对于$f\in\mathrm{Lip}\,\alpha$，
\[
\|\sigma_nf-f\|_\infty=O(n^{-\alpha/2}).
\]
与傅里叶部分和相比，Fejér算子彻底消除了对数因子，但也付出了代价——它的逼近阶在光滑性较高时只有$n^{-\alpha/2}$，而非$S_n$的$n^{-\alpha}$（不计对数项）。这看似矛盾：为什么范数更小的算子，逼近速度反而（在$\alpha$较大时）更慢？答案在于$S_n$是投影，它能复制$\mathcal{T}_n$中的任何多项式，因此在光滑函数上可以充分利用函数的“高阶”信息；而$\sigma_n$丧失了投影性，本质上只利用了一阶光滑性。这一现象称为线性正算子的\textbf{饱和}（saturation），我们将在下一小节详细讨论。

\subsection{Zygmund类与逼近阶的充要条件}

为了从“函数光滑性”和“逼近阶”两个方向同时刻画，Zygmund引入了一类比Lipschitz条件更精细的函数类。对$0<\alpha<1$，$\mathrm{Lip}\,\alpha$已经足够；但对$\alpha=1$，Lipschitz条件本身并不足以给出$E_n(f)=O(1/n)$的充要特征。Zygmund类$\Lambda_*$（或记作$Z$）定义为满足
\[
|f(x+h)+f(x-h)-2f(x)|\le Mh
\]
的连续函数。这比Lipschitz条件弱，但它恰是$E_n(f)=O(1/n)$的充要条件。

更一般地，可定义Zygmund类$Z^\alpha$（$\alpha>0$），其与最佳逼近有如下精确对应：
\[
f\in Z^\alpha\iff E_n(f)=O(n^{-\alpha}).
\]
这就是著名的\textbf{Jackson–Zygmund定理}（正定理与逆定理的结合）。你课上学过的例子：$|x|$和$|\sin x|$都属于$Z^1$，因此它们的最佳三角多项式逼近阶恰好为$O(1/n)$。注意这里的最佳逼近是$E_n(f)$，而不是某个具体算子的逼近误差。

回到我们的两个算子上。由Lebesgue不等式或直接估计，可得
\[
\|f-S_n(f)\|_\infty\le (1+\Lambda_n)E_n(f),\qquad
\|\sigma_nf-f\|_\infty\le C\omega(f,n^{-1/2}).
\]
对于$f\in Z^1$（比如$|x|$或$|\sin x|$），$E_n(f)=O(1/n)$。因此
\[
\|f-S_n(f)\|_\infty=O\!\left(\frac{\ln n}{n}\right),\qquad
\|\sigma_nf-f\|_\infty=O(n^{-1/2}).
\]
傅里叶部分和虽然带有对数因子，但其阶$n^{-1}\ln n$远优于Fejér的$n^{-1/2}$。这说明，在处理具有较好光滑性的函数（特别是二阶广义导数存在的函数）时，投影型算子在逼近阶上具有显著优势。但投影性质的代价是Lebesgue常数发散，导致它在$C_{2\pi}$上不是一致有界的，从而不保证对所有连续函数收敛。

Fejér算子用丧失投影性质的代价换取了全空间的一致收敛性和最稳定的范数（$1$），但它在光滑函数上会遭遇饱和：即使$f$无限可微，$\sigma_nf-f$的衰减速度也无法超越$O(1/n)$（实际上，对于Fejér算子，饱和阶是$O(1/n)$，此处的$O(n^{-1/2})$是对于Lipschitz类的估计；关于饱和的详细讨论可参考DeVore \& Lorentz）。这正是线性正算子的本质局限。

\subsection{小结：Korovkin–Lebesgue–Jackson框架下的统一图景}

我们现在可以将一切串联起来：
\begin{itemize}
\item \textbf{Korovkin定理}提供定性收敛判据：只需三个测试函数，保证线性正算子序列在整个空间上一致收敛。
\item \textbf{Lebesgue不等式}（或其定量替代）将构造性算子的误差与最佳逼近误差挂钩，其中的Lebesgue常数衡量了算子的稳定性。
\item \textbf{Jackson定理}及\textbf{Zygmund类}给出最佳逼近误差随光滑性的衰减律，从而导出各算子在具体函数类上的逼近阶。
\end{itemize}

对于傅里叶部分和，其Lebesgue常数对数发散，导致逼近阶中常含有对数因子；对于Fejér算子，常数恒为$1$，但因其非投影性，逼近阶受连续模控制，在光滑函数上出现饱和，仅为$O(n^{-1/2})$（或$O(1/n)$饱和）。两者互补：一个在光滑类上逼近更快但不稳定，一个稳定但速率较慢。这深刻揭示了逼近论中“最优”与“可构造”之间永恒的张力。
\end{document}